\section{Problem Formulation}
\label{sec:formulation}

\subsection{Notation}
Let $\mathcal{K}=\{1,\dots,K\}$ denote the set of FL clients.
Each client $k$ holds local benign traffic data $\mathcal{D}_k^{\text{ben}}$ and may have
attack data $\mathcal{D}_k^{\text{atk}}$.
An autoencoder $h = g \circ f$ maps input $\mathbf{x} \in \mathbb{R}^d$ to a reconstruction
$\hat{\mathbf{x}} = g(f(\mathbf{x}))$, where $f$ is the encoder and $g$ the decoder.
The reconstruction error is the per-sample mean squared error:
\begin{equation}
  e(\mathbf{x}) = \frac{1}{d}\|\mathbf{x} - \hat{\mathbf{x}}\|_2^2,
\end{equation}
where $d$ is the input dimension.

\subsection{Binary Alarm Decision}
Given a threshold $\tau_k$ for client $k$, the alarm decision is:
\begin{equation}
  \text{anomaly}(\mathbf{x}) = \mathbf{1}[e(\mathbf{x}) > \tau_k].
\end{equation}
The FPR for client $k$ is $\rho_k = \Pr[e(\mathbf{x}) > \tau_k \mid \mathbf{x} \sim \mathcal{D}_k^{\text{ben}}]$.

\subsection{Dispersion Metric}
The primary metric is the coefficient of variation of FPR over eligible clients
$\mathcal{K}_{\mathrm{elig}} \subseteq \mathcal{K}$:
\begin{equation}
  \text{CV(FPR)} = \frac{\sigma_{\rho}}{\mu_{\rho}},\quad
  \mu_{\rho} = \frac{1}{|\mathcal{K}_{\mathrm{elig}}|}\sum_{k\in\mathcal{K}_{\mathrm{elig}}} \rho_k,
\end{equation}
where $\sigma_{\rho}$ is the standard deviation.
CV(FPR) normalizes cross-client FPR dispersion by the fleet mean, making dispersion
comparable across regimes with different mean FPR levels; we report IQR(FPR) and
max--min FPR as absolute-dispersion checks to guard against small-denominator artifacts.

\textbf{Eligibility and fallback.}
Eligible clients $\mathcal{K}_{\mathrm{elig}}$ are those with at least $n_{\min}{=}100$ benign calibration samples.
Clients below this threshold are \emph{calibration-pending}: they receive the client-averaged
shared threshold $\tau_{\text{global}}$ (B1 value) as a fallback and are excluded from
dispersion metrics.
Coverage is reported as $|\mathcal{K}_{\mathrm{elig}}|/|\mathcal{K}|$ for all conditions.

\textbf{Operational assumption.}
Percentile calibration assumes each client's benign reconstruction-error distribution
is sufficiently stable between calibration and test.
Temporal drift would weaken the calibration guarantee; this paper does not model concept drift.

\subsection{Confirmatory Statistic}
The primary statistical estimate is the seed-level delta:
\begin{equation}
  \Delta_s = \text{CV(FPR)}_{\text{B1},s} - \text{CV(FPR)}_{\text{B2},s},
\end{equation}
aggregated over seeds $s \in \mathcal{S}$ via descriptive summaries: mean, standard deviation, range, and sign consistency.
With only five seeds, we report these statistics as a seed-level consistency summary---not a high-powered
population significance test.
We interpret all-positive seed deltas as seed-level support
for RQ2 in this tested protocol.

\subsection{Research Questions}
\begin{itemize}
  \item \textbf{RQ1:} Under the client-averaged shared threshold (B1), how large is the variation of FPRs across heterogeneous IoT clients?
  \item \textbf{RQ2:} Does per-client threshold calibration (B2) reduce CV(FPR) compared with the client-averaged shared threshold, and is the reduction consistent across seeds?
  \item \textbf{RQ3:} How much of the B2 dispersion reduction can cluster-mean calibration (B4) recover?
\end{itemize}

The benign calibration-error distributions of different device types differ substantially
in location and spread, so the B1 client-averaged threshold $\tau_{\text{global}}$ crosses
each distribution at a different fraction, producing device-specific FPRs.
Per-client percentile calibration is expected to align benign-tail exceedance rates
when calibration and test distributions are sufficiently stable; the empirical questions
are therefore the magnitude of the FPR-dispersion reduction, the detection-quality
tradeoff it creates, and the extent to which grouped calibration (B4) approximates
per-client calibration without a device taxonomy---all under a fixed FL-AE setting
with real physical-device heterogeneity.
